\scanpage{112}\mspage{6}
Montrons comment cela permet de calculer $z$, via $z_0=p(z)$ et $z'=p'(z)$ (reliés par
\[
 {z'}^2=z_0^3+a_1z_0^2+a_2z_0+a_3\qquad), \tag*{(46)}
\]
quand on connaît $x$ et $y$, exprimés à l'aide de \deleted{$x_0=p(x),x'=p'(x)$}
\[
 \left\{\begin{aligned}
 x_0&=p(x),\ x'=p'(x)\\
 y_0&=p(y),\ y'\mathrel{\unsure{-|=}}p'\unsure{y)|(y)}
 \end{aligned}\right.
\]
(reliés par
\[
 {x'}^2=x_0^3+a_1x_0^2+a_2x_0+a_3,\quad
 {y'}^2=y_0^3+a_1y_0^2+a_2y_0+a_3).
\]
L'équation (\unsure{41|4)|42}) s'écrit
\[
 z'(x_0-y_0)+z_0(y'-x')+(x'y_0-y'x_0)=0 \tag*{$(*)$}
\]
i.e.
\[
 z'=\lambda z_0+\mu,\quad\text{avec}\quad
 \lambda=\frac{x'-y'}{x_0-y_0},\quad
 \mu=\frac{y'x'_0-x'y_0}{x_0-y_0} \tag*{(47)}
\]
\deleted{Donc} (NB on a $x_0\ne y_0$ i.e $x_0-y_0\ne0$).

De plus $z'$ satisfait (46), qui équivaut \deleted{à} \deleted{\unsure{l'équa|l'équation}} (compte tenu de (47), qui implique $ {z'}^2=\lambda^2z_0^2+2\lambda\mu z_0+\mu^2$) à l'équation
\[
 z_0^3+(a_1-\lambda^2)z_0^2+(a_2-2\lambda\mu)z_0+(a_3-\mu^2)=0. \tag*{(48)}
\]
Donc $(z_0,z')$ est lié par (47), (48), a priori on trouve 3 solutions \inserted{en $z_0$} de (48), à qui correspondent trois solutions $(z_0,z')$ via (47). Il faut trouver une solution unique, on se rappelle qu'on doit avoir
\[
 z_0\ne x_0,\quad z_0\ne y_0,\qquad\deleted{\text{\unsure{do|de}}} \tag*{(49)}
\]
% archive sheet number: 111
